% Preámbulo común del material docente · XeLaTeX
\documentclass[11pt,a4paper]{article}
\usepackage[margin=22mm,top=20mm,bottom=22mm]{geometry}
\usepackage{fontspec}
\setmainfont{Avenir Next}[BoldFont={Avenir Next Demi Bold}, BoldItalicFont={Avenir Next Demi Bold Italic}]
\setsansfont{Avenir Next}
\setmonofont{Menlo}[Scale=0.88]
\newfontfamily\symfont{Apple Symbols}
\newfontfamily\unifont{Arial Unicode MS}
\usepackage{unicode-math}
\setmathfont{latinmodern-math.otf}
\usepackage{polyglossia}\setdefaultlanguage{spanish}
\usepackage{xcolor,booktabs,tabularx,enumitem,parskip,microtype,titlesec,fancyhdr}
\usepackage[most]{tcolorbox}
\usepackage[hidelinks]{hyperref}
\emergencystretch=2em

\definecolor{tinta}{HTML}{1F2933}
\definecolor{acento}{HTML}{B7410E}
\definecolor{suave}{HTML}{F4F1EC}
\definecolor{gris}{HTML}{6B7280}
\definecolor{linea}{HTML}{D6D0C6}
\color{tinta}

\titleformat{\section}{\Large\bfseries\color{acento}}{\thesection}{0.6em}{}
\titleformat{\subsection}{\large\bfseries}{\thesubsection}{0.6em}{}
\titlespacing*{\section}{0pt}{1.6em}{0.5em}
\titlespacing*{\subsection}{0pt}{1.1em}{0.3em}
\setlist{nosep, leftmargin=1.4em, itemsep=2pt}
\renewcommand{\arraystretch}{1.25}

% bloque de órdenes de terminal
\newtcblisting{orden}{listing only, colback=suave, colframe=linea, boxrule=0.4pt, arc=2pt,
  left=6pt, right=6pt, top=4pt, bottom=4pt, listing options={basicstyle=\ttfamily\small, breaklines=true, columns=fullflexible, keepspaces=true}}
% nota destacada
\newtcolorbox{nota}[1][Nota]{colback=white, colframe=acento, boxrule=0.6pt, arc=2pt, title=#1,
  fonttitle=\bfseries\small, coltitle=white, colbacktitle=acento, left=6pt, right=6pt}
% preguntas
\newcommand{\pregunta}[2]{\item[\textbf{(#1)}] #2}
\newenvironment{preguntas}{\begin{description}[leftmargin=2.4em, style=nextline, itemsep=4pt]}{\end{description}}
\newcommand{\ok}{{\symfont ✔}}
\newcommand{\nok}{{\symfont ✘}}
\newcommand{\qed}{{\symfont ∎}}
\newcommand{\codigo}[1]{\texttt{#1}}

\pagestyle{fancy}\fancyhf{}
\renewcommand{\headrulewidth}{0pt}
\fancyfoot[C]{\small\color{gris}\docpie\ · página \thepage}

\newcommand{\cabecera}[3]{%
  {\Huge\bfseries\color{acento}#1\par}\vspace{4pt}
  {\large #2\par}\vspace{2pt}
  {\color{gris}#3\par}\vspace{8pt}
  {\color{linea}\hrule height 0.6pt}\vspace{10pt}}

\newcommand{\docpie}{Laboratorio 1 · notas del docente}
\begin{document}
\cabecera{Laboratorio 1 · notas del docente}{Bisección con un agente de RL · 90 minutos}{Resultados medidos el 22 de septiembre de 2026 en un MacBook con Python 3.13, semilla 1. Todo se reproduce con \codigo{docs/curso/lab01/experimentos.py}.}

\section{Tiempos}

\begin{tabularx}{\linewidth}{@{}X l@{}}
\toprule
\textbf{comando} & \textbf{duración} \\
\midrule
\codigo{python -m bisectrl} (con animación, 3000+3000 gen) & 1–2 min; con \codigo{--speed 3} unos 40 s \\
\codigo{python -m bisectrl --no-tui} & 0,35 s \\
\codigo{experimentos.py generaciones} & 1 s \\
\codigo{experimentos.py recompensa} & 1 s \\
\codigo{experimentos.py sin\_medio} / \codigo{grafo} & $<1$ s \\
\bottomrule
\end{tabularx}

La animación necesita una terminal de $120\times 40$. Si los monitores del laboratorio son pequeños, proyectar una sola corrida y que las parejas trabajen con \codigo{--no-tui} y el informe.

\section{Plan de la sesión}

\begin{tabularx}{\linewidth}{@{}l X l@{}}
\toprule
\textbf{min} & \textbf{bloque} & \textbf{modo} \\
\midrule
0–30 & terminal desde cero: abrir, \codigo{cd}, \codigo{ls}, descargar ZIP, \codigo{venv}, \codigo{pip install}. Ir puesto por puesto. Meta: todos corren \codigo{--no-tui} y ven un informe & guiado \\
30–40 & RL en la pizarra con la bisección como único ejemplo: estado, acción, recompensa, política. Preguntas (a)–(c) en papel & expositivo \\
40–55 & parte 2: correr con animación, anotar hitos, leer la síntesis & parejas \\
55–70 & parte 3: romper la recompensa; plenaria breve sobre (f) & parejas + plenaria \\
70–85 & parte 4: dibujar el grafo, experimento \codigo{grafo}, plenaria sobre (n) & parejas + plenaria \\
85–90 & cierre y entrega & expositivo \\
\bottomrule
\end{tabularx}

Las dos ideas de la sesión son (f) y (n). Si se va el tiempo, sacrificar \codigo{sin\_medio} y la comparación de semillas, nunca esas dos plenarias.

\begin{nota}[Riesgo principal: la instalación]
Probar el día anterior en un PC del laboratorio que \codigo{python --version} funciona, que hay internet para \codigo{pip}, y que Windows Terminal se abre. Si \codigo{pip} no tiene red, llevar un USB con \codigo{rich} y \codigo{pytest} descargados (\codigo{pip download rich pytest}) e instalar con \codigo{--no-index --find-links}.
\end{nota}

\section{Resultados esperados y respuestas}

\begin{preguntas}
\pregunta{a}{Con $\lambda=0{,}5$ el ancho se divide por 2: $-1+\log_2 2 = 0$ puntos por paso, exacto. Con $\lambda=0{,}1$, el 10\,\% de las veces (aproximadamente, si la raíz cae uniforme) el progreso es $\log_2 10\approx 3{,}3$ y el 90\,\% es $\log_2(1/0{,}9)\approx 0{,}15$; en promedio unos $0{,}47$ bits, o sea $Q(0{,}1)\approx -0{,}53$. Medido: $-0{,}55$ a $-0{,}76$ según semilla. La raíz no cae uniforme, así que las cifras varían.}
\pregunta{b}{Sí: los tres signos bastan. Es el punto que conecta con la demostración: el invariante $f(a_n)\,f(b_n)\le 0$ es exactamente lo que la regla de conservación mantiene.}
\pregunta{c}{$\lambda=1/2$. Argumento minimax: el peor caso de cualquier $\lambda\neq 1/2$ es peor que el de $1/2$. Los mejores estudiantes notarán que el promedio también lo favorece si la raíz está distribuida simétricamente.}
\end{preguntas}

\textbf{Hitos con semilla 1} (3000+3000 generaciones):

\begin{tabularx}{\linewidth}{@{}X r@{}}
\toprule
\textbf{hito} & \textbf{generación} \\
\midrule
primera convergencia & 7 \\
$\lambda=1/2$ estable (300 gen seguidas) & 345 \\
invariante de Bolzano & 358 \\
$\ge 95\,\%$ de convergencia & 391 \\
primera demostración completa & 256 (fase 2) \\
política voraz produce la mínima & 300 \\
primera demostración mínima durante el entrenamiento & 459 \\
\bottomrule
\end{tabularx}

Con semilla 42 y solo 200+100 generaciones (una corrida vieja en \codigo{runs/}) no aprende nada: es un buen ejemplo de «el informe dice honestamente que no lo logró».

\begin{preguntas}
\pregunta{d}{$Q(\lambda)$ es el progreso medio por evaluación menos 1. $Q(0{,}5)=0$ porque el progreso es exactamente 1 bit. Valores negativos = menos de 1 bit por evaluación en promedio.}
\pregunta{e}{Las generaciones cambian (decenas arriba o abajo); lo aprendido, no. Es el punto para distinguir \emph{proceso estocástico} de \emph{resultado}.}
\pregunta{f}{Resultado medido con «perder la raíz cuesta 0»: \textbf{0 convergencias, 2000 raíces perdidas}, éxito 0\,\%. El agente aprende a tirar la raíz en el primer paso. Es racional: un episodio que pierde la raíz de inmediato gana $-1+\log_2(1/(1-\lambda))$ una sola vez y termina; bisecar 12 veces gana 0 en total (con $\lambda=1/2$ cada paso vale exactamente 0). Terminar rápido \emph{domina}. Es el ejemplo más limpio de \emph{reward hacking} que se puede dar a este nivel, y conecta con la idea de que la recompensa es la especificación completa del problema.}
\pregunta{g}{Con paso gratis el agente sigue eligiendo $\lambda=1/2$. La cabeza de corte usa $\gamma=0$ y tasa $1/N$: maximiza el progreso esperado \emph{de este paso}, y ese progreso lo maximiza el punto medio con o sin el $-1$. Buen momento para decir que la suma de las recompensas de progreso telescopa: $\sum\log_2(w_{k-1}/w_k)=\log_2(w_0/w_n)$, así que el shaping no cambia qué es óptimo (Ng, Harada y Russell, 1999). No hace falta citar el artículo a los estudiantes.}
\pregunta{h}{Respuestas típicas buenas: «premiar por cada evaluación de $f$» (evaluaría infinito), «premiar que $|f(x)|$ sea pequeño» (se estanca donde $f$ es plana), «premiar solo al final» (aprende más lento pero no se puede engañar). Cualquier respuesta que identifique \emph{qué haría el agente} vale.}
\pregunta{i}{Medido sin $0{,}5$: elige $\lambda=0{,}4$ con $Q(0{,}4)=-0{,}029$, $Q(0{,}6)=-0{,}039$; la curva es simétrica en forma de «V invertida» con máximo en $0{,}5$. Las diferencias $0{,}4/0{,}6$ son ruido de muestreo; pedir que corran otra semilla para verlo.}
\pregunta{j}{El grafo tiene profundidad 7 (H $\to$ DEF $\to$ MONO $\to$ CONV $\to$ SAME $\to$ CONT $\to$ ROOT $\to$ QED). Órdenes válidos: muchos; por ejemplo H1, H2, DEF, INV, WIDTH, MONO\_A, MONO\_B, CONV\_A, CONV\_B, SAME, CONT, ROOT, MID, QED. Basta con que respeten las flechas.}
\pregunta{k}{\codigo{D\_BOLZ} es el importante: existencia de una raíz (Bolzano) no dice nada sobre a dónde converge la sucesión de puntos medios; puede haber varias raíces y hay que probar que $m_n$ converge y que su límite es raíz. Otros: \codigo{D\_LIP}, Lipschitz no está en las hipótesis y «luego es de Cauchy» no se sigue; \codigo{D\_UNIQ}, falso en general; \codigo{D\_TVM}, verdadero bajo derivabilidad pero irrelevante; \codigo{D\_DERIV}, \codigo{D\_NEWTON}, otro método; \codigo{D\_DIV}, falso (acotada).}
\pregunta{l}{El agente aprendió un orden topológico del grafo y a evitar nodos sin salida. No «entendió» nada. Esta es la respuesta que hay que arrancarles; si dicen «el agente demostró el teorema», volver al experimento \codigo{grafo}.}
\pregunta{m}{\codigo{ROOT} dice: como $f(a_n)\,f(b_n)\le 0$ para todo $n$ y el límite es $f(c)^2$, entonces $f(c)^2\le 0$. Sin \codigo{INV} no sabemos que los productos son $\le 0$; el límite de $f(a_n)\,f(b_n)$ sigue siendo $f(c)^2$, pero nada obliga a que sea $\le 0$. El verificador acepta la demostración de 13 pasos porque nadie le dijo que faltaba algo. Medido: el verificador devuelve \codigo{finished = True} para el orden sin \codigo{INV}.}
\pregunta{n}{El grafo es la demostración. Quien escribe el grafo demuestra; el agente lo ordena. Esto justifica el proyecto final: escribir un grafo correcto es el trabajo matemático.}
\pregunta{tarea}{Un distractor es «tentador» para el agente solo por el bono de exploración (UCB): lo intenta un número parecido de veces a cualquier otro nodo inválido hasta que su $Q$ cae. Para un humano es tentador si \emph{parece} seguirse. El agente no tiene noción de plausibilidad; el humano sí. En el informe, la tabla «saltos lógicos más frecuentes» muestra los intentos por paso.}
\end{preguntas}

\textbf{Extras si preguntan.} Estados de la regla de conservación: 3 signos con $f(x)$ posiblemente 0 dan $2\cdot 3\cdot 2=12$ combinaciones, pero con el invariante $f(a)$ y $f(b)$ tienen signos opuestos, así que solo 6 son alcanzables (4 si $f(x)\neq 0$). Número de generaciones (\codigo{experimentos.py generaciones}, no está en la guía de 90 min): con 100 el éxito es 15\,\% y no hay demostración completa; con 300, 74\,\% y la política voraz ya produce la mínima aunque durante el entrenamiento no salió ninguna ($\varepsilon$ todavía alto); con 1000 ya está todo.

\section{Trampas de la sesión}
\begin{itemize}
\item Si editan \codigo{proof\_kb.py} y olvidan revertir, la siguiente sesión arranca con el grafo alterado. Insistir en \codigo{git checkout bisectrl/proof\_kb.py} (o en volver a descomprimir el ZIP).
\item Si un distractor nuevo tiene \codigo{deps} no vacías o \codigo{distractor=False}, cambia \codigo{REQUIRED\_KEYS}. Solo hace falta \codigo{Step("D\_MIO", "Distractor", "...", distractor=True)}.
\item \codigo{keep\_policy\_is\_bolzano()} devuelve \codigo{True} con solo 4 estados visitados y coherentes; con 100 generaciones puede dar \codigo{True} aunque el éxito sea 15\,\%. No es contradicción: la regla ya es correcta pero el corte todavía es malo. Buena pregunta extra.
\item Terminal pequeña: \codigo{rich} recorta paneles y parece que «falta» información. Usar \codigo{--no-tui}.
\end{itemize}

\section{Qué se evalúa}

Rúbrica sugerida (10 puntos): haber corrido el agente y entregado la tabla de hitos (2), (a)–(e) (2), (f)–(i) (3), (j)–(n) (3). Pesar (f) y (n): son las dos ideas de la sesión. El distractor propio es opcional y suma hasta 1 punto extra.
\end{document}
